\section*{الفصل \ref{ch:b3:trace-element-modelling} --- نمذجة العناصر النزرة}

\begin{solution}{exo:b3:trace-element-modelling:1}
الليرزوليت السبينيلي: $0.12\times0.40 + 0.88\times0.016 = 0.048 + 0.0141 =
0.062$. الليرزوليت الجارنيتي: $0.195\times0.40 + 0.167\times0.11 +
0.639\times0.016 = 0.078 + 0.0184 + 0.0102 = 0.107$.
\end{solution}

\begin{solution}{exo:b3:trace-element-modelling:2}
$1/(0.05 + 0.02\times0.95) = 1/0.069 = \qty{14.5}{ppm}$;
$1/(0.05 + 0.10\times0.95) = 1/0.145 = \qty{6.9}{ppm}$.
\end{solution}

\begin{solution}{exo:b3:trace-element-modelling:3}
يبلغ الإثراء أقصاه حين $F \to 0$: $1/D = 100$. ومع $D = 0$،
$C_L/C_0 = 1/F = 20$.
\end{solution}

\begin{solution}{exo:b3:trace-element-modelling:4}
$\mathrm{La/Yb} = 30/2.5 = 12.0$; $(\mathrm{La}/\mathrm{Yb})_N =
(30/0.232)/(2.5/0.163) = 129.3/15.34 = 8.4$.
\end{solution}

\begin{solution}{exo:b3:trace-element-modelling:5}
\omterm{thm:b3:trace-element-modelling:batch}{الانصهار الدفعي}: $C_S/C_0 = 0.01/(0.01 + 0.05\times0.99) = 0.168$. \omterm{thm:b3:trace-element-modelling:fractional}{الانصهار التجزيئي}:
$(0.95)^{99} = 0.0062$، أي أقلّ بمقدار 27 مرة. ويحتوي المصهور اللحظي
$0.0062/0.01 = 0.62\,C_0$: أقلّ من المصدر.
\end{solution}

\begin{solution}{exo:b3:trace-element-modelling:6}
$D_0 = 0.58\times0.0003 + 0.27\times0.002 + 0.12\times0.053 + 0.03\times0.01
= 0.00737$; $P = 0.20\times0.0003 + 0.25\times0.002 + 0.45\times0.053 +
0.10\times0.01 = 0.0254$. عند $F = 0.02$: \omterm{def:b3:trace-element-modelling:nonmodal}{الانصهار غير الشكلي} $1/(0.00737 +
0.02\times0.9746) = 37.2$، والانصهار الشكلي $1/(0.00737 + 0.02\times0.9926) = 36.7$. وعند
$F = 0.10$: 9.54 و9.38، بفارق \qty{1.7}{\percent}.
\end{solution}

\begin{solution}{exo:b3:trace-element-modelling:7}
$D = 0.1$: \omterm{thm:b3:trace-element-modelling:batch}{الانصهار الدفعي} $1/(0.1 + 0.09) = 5.26$؛ \omterm{def:b3:trace-element-modelling:fractional}{والمصهور التجزيئي المجمَّع} $(1 - 0.9^{10})/0.1 =
6.51$. $D = 0.01$: \omterm{thm:b3:trace-element-modelling:batch}{الانصهار الدفعي} $1/(0.01 + 0.099) = 9.17$؛ \omterm{def:b3:trace-element-modelling:fractional}{والمصهور التجزيئي المجمَّع}
$(1 - 0.9^{100})/0.1 = 10.00$. يختلف الزوج ذو $D = 0.1$ أكثر: أمّا
$D = 0.01$ فيحتوي كلا المصهورين أصلًا على جلّ العنصر تقريبًا، ويميل كلاهما إلى
$1/F$.
\end{solution}

\begin{solution}{exo:b3:trace-element-modelling:8}
$D' = 0.99\times0.001 + 0.01 = 0.01099$. المصهور الأول: $1/D' = 91$، مقابل
1000 في \omterm{thm:b3:trace-element-modelling:fractional}{الانصهار التجزيئي}. عند $F = 0.01$: $91\times0.99^{90.0} = 36.8$،
مقابل $1000\times0.99^{999} = 0.044$: يحفظ المصهور المسامي الباقي من أن يُستنزف.
أن يُستنزف.
\end{solution}

\begin{solution}{exo:b3:trace-element-modelling:9}
ترتفع النسبة بعامل قدره $(D_{\mathrm{Sm}} + F(1 - D_{\mathrm{Sm}}))/(D_{\mathrm{La}}
+ F(1 - D_{\mathrm{La}}))$: عند $F = 0.01$، $0.0714/0.0177 = 4.0$؛ وعند $F = 0.10$،
$0.156/0.107 = 1.46$. تكون النسبة حسّاسة للقيمة $F$ عندما تكون $F$ صغيرة، وتكاد
تكون عديمة الحساسية فوق نحو \qty{10}{\percent}؛ وتعتمد كذلك على $D_{\mathrm{Sm}}$،
ومن ثمّ على المعادن الباقية.
\end{solution}

\begin{solution}{exo:b3:trace-element-modelling:10}
$\frac1F\int_0^F \frac1D(1 - x)^{1/D - 1}\,\mathrm dx = \frac1F\bigl[-(1 -
x)^{1/D}\bigr]_0^F = [1 - (1 - F)^{1/D}]/F$، الصيغة الواردة في الفصل، إذا
$(1 - F)^{1/D} \to 0$ فإنّها تؤول إلى $1/F$ مضروبًا في $C_0$. ضمن \qty{1}{\percent}:
$(1 - F)^{100} < 0.01$، $F > 1 - 0.01^{0.01} = 0.045$، أي نحو
\qty{4.5}{\percent}.
\end{solution}

\begin{solution}{exo:b3:trace-element-modelling:11}
الميل $a = (10 - 6)/(20 - 10) = \qty{0.40}{ppm^{-1}}$؛ والمقطع $b = 10 -
0.40\times20 = 2.0$. $D_{\mathrm{Yb}} = 0.40\times0.44 = 0.176$;
$C_0^{\mathrm{La}} = 2.0\times0.44/0.824 = \qty{1.07}{ppm}$. $F_P = 1.07/20 =
\qty{5.3}{\percent}$, $F_Q = 1.07/10 = \qty{10.7}{\percent}$.
\end{solution}

\begin{solution}{exo:b3:trace-element-modelling:12}
$F = x_i/p_i$: الكلينوبيروكسين $0.12/0.45 = 0.27$، والسبينل $0.03/0.10 = 0.30$،
الأورثوبيروكسين $0.27/0.25 = 1.08$، والأوليفين $0.58/0.20 = 2.9$؛ ينفد الكلينوبيروكسين
أوّلًا. عند $F = 0.2$: $\bar D = (0.00737 - 0.0254\times0.2)/0.8 =
0.0029$. وبعد $F = 0.27$ تتغيّر القيم $p_i$ (إذ لم يبقَ كلينوبيروكسين لينصهر)،
ولم تعد الصيغة، التي تفترض ثبات $p_i$، تنطبق.
\end{solution}

\begin{solution}{pb:b3:trace-element-modelling:1}
\textbf{1.} $D_{\mathrm{La}} = 0.12\times0.058 + 0.88\times0.001 = 0.00784$;
$D_{\mathrm{Yb}} = 0.12\times0.58 + 0.88\times0.16 = 0.210$.

\textbf{2.} $D_{\mathrm{La}} = 0.195\times0.058 + 0.167\times0.001 +
0.639\times0.001 = 0.0121$; $D_{\mathrm{Yb}} = 0.195\times0.58 +
0.167\times4.5 + 0.639\times0.16 = 0.967$. بعد التقريب: 0.0078، 0.21، 0.012، 0.97.

\textbf{3.} $D_{\mathrm{Yb}}$، من 0.21 إلى 0.97، لأنّ الجارنيت يحتوي على
العناصر الأرضية النادرة الثقيلة ($D = 4.5$). ويبقى الإتربيوم غير متوافق بدرجة طفيفة جدًّا،
$0.97 < 1$.

\textbf{4.} $(0.648/0.232)/(0.441/0.163) = 2.79/2.71 = 1.03$: نمط مسطّح
نحو 2.8 مرة قيمة الكوندريت.

\textbf{5.} $0.648/(0.0078 + 0.01\times0.9922) = \qty{36.6}{ppm}$؛ عند
\qty{5}{\percent}، \qty{11.3}{ppm}؛ وعند \qty{10}{\percent}، \qty{6.05}{ppm}.

\textbf{6.} $F = (0.648/9.6 - 0.0078)/0.9922 = (0.0675 - 0.0078)/0.9922 =
0.0602$, \qty{6.0}{\percent}.

\textbf{7.} $\mathrm{Yb} = 0.441/(0.21 + 0.0602\times0.79) = \qty{1.71}{ppm}$
و\mbox{\textdir TLT La/Yb $= 5.6$}؛ الحمم A لها 1.7 و5.6، فهي تتفق مع مصدر سبينلي.

\textbf{8.} \mbox{\textdir TLT La $= 0.0078\times9.6 = \qty{0.075}{ppm}$}، \mbox{\textdir TLT Yb} $= 0.21\times1.71 =
\qty{0.36}{ppm}$. اللانثانوم المتبقّي: $0.940\times0.075/0.648 = 0.11$، أي نحو
\qty{11}{\percent}.

\textbf{9.} $F = 0.648/9.6 = \qty{6.75}{\percent}$.

\textbf{10.} $(0.648/0.0078)(1 - F)^{127.2}$: \qty{73}{ppm} عند
\qty{0.1}{\percent}، \qty{23}{ppm} عند \qty{1}{\percent}، \qty{0.03}{ppm} عند
\qty{6}{\percent}.

\textbf{11.} $0.648\times(1 - 0.940^{128.2})/0.060 = \qty{10.8}{ppm}$.

\textbf{12.} $F = \qty{6.75}{\percent}$: عند هذا $F$، $(1 - F)^{1/D} =
\num{1.3e-4}$، فيحتفظ المصهور المجمَّع بكامل اللانثانوم الموجود في المصدر،
و$\bar C_L = C_0/F$ كما في عنصر غير متوافق تمامًا.

\textbf{13.} $0.648\times0.9325^{127.2} = \qty{9e-5}{ppm}$، أي أقلّ بنحو 800 مرّة
من \qty{0.075}{ppm} في \omterm{thm:b3:trace-element-modelling:batch}{الانصهار الدفعي}.

\textbf{14.} يحصر النموذجان الجواب بين \qty{6.0}{\percent} (\omterm{thm:b3:trace-element-modelling:batch}{الانصهار الدفعي})
و\qty{6.75}{\percent} (\omterm{def:b3:trace-element-modelling:fractional}{المصهور التجزيئي المجمَّع}). يحدّد اللانثانوم $F$ ضمن
نحو \qty{1}{\percent} من الانصهار، أيًّا كان النموذج.

\textbf{15.} الحمم A: \mbox{\textdir TLT La/Yb $= 5.6$}، $(\mathrm{La}/\mathrm{Yb})_N =
(9.6/0.232)/(1.7/0.163) = 4.0$. الحمم B: \mbox{\textdir TLT La/Yb $= 34.7$}، $(\mathrm{La}/\mathrm{Yb})_N
= (15.6/0.232)/(0.45/0.163) = 24.4$.

\textbf{16.} عندما $F \to 0$: $(0.648/0.0078)/(0.441/0.21) = 83.1/2.10 = 39.6$.

\textbf{17.} $F = \qty{0.11}{\percent}$، حيث يحتوي المصهور على
\qty{72}{ppm} من La و\qty{2.1}{ppm} من Yb. الحمم B تحتوي على \qty{15.6}{ppm} من La:
لا يمكن لمصدر سبينلي أن ينتجها.

\textbf{18.} $F = (0.648/15.6 - 0.012)/0.988 = 0.0299$، \qty{3.0}{\percent}؛
$\mathrm{Yb} = 0.441/(0.97 + 0.0299\times0.03) = \qty{0.454}{ppm}$، كما
لوحظ.

\textbf{19.} يحلّ الجارنيت محلّ السبينل على عمق نحو 60 إلى \qty{90}{km}: بدأت الحمم B
تنصهر على عمق أكبر من ذلك، والحمم A على عمق أقلّ، في نطاق السبينل. ومصهور الليرزوليت الجارنيتي
بمقدار \qty{9.6}{ppm} من La ($F \approx \qty{5.6}{\percent}$) سيكون له La/Yb
قريبًا من 21، لا 5.6: الحمم A لم تأتِ من الليرزوليت الجارنيتي.

\textbf{20.} $a = 0.97/0.441 = \qty{2.20}{ppm^{-1}}$, $b = 0.648\times0.03/0.441
= 0.044$؛ عند \mbox{\textdir TLT La $= 15.6$}، تكون \mbox{\textdir TLT La/Yb $= 2.20\times15.6 + 0.044 = 34.4$}، مقابل
34.7 للحمم B.

\textbf{21.} يرفض الأوليفين العنصرين كليهما: نسبة La/Yb لا تتغيّر. المصهور الأوّلي
كان يحتوي على $9.6\times0.9 = \qty{8.64}{ppm}$ من La، لذا $F = (0.648/8.64 -
0.0078)/0.9922 = \qty{6.8}{\percent}$ بدلًا من 6.0.

\textbf{22.} يمكن لمصدر واحد أن ينتج الحممين معًا: الحمم B مصهور صغير
(\qty{3}{\percent}) بدأ في نطاق الجارنيت، على عمق يزيد عن نحو 60 إلى
\qty{90}{km}؛ والحمم A مصهور أكبر في نطاق السبينل.

\textbf{23.} الحمم A ناتجة عن \omterm{thm:b3:trace-element-modelling:batch}{الانصهار الدفعي} بنسبة \textbf{\qty[text-series-to-math=true]{6.0}{\percent}}
من ليرزوليت سبينلي (\qty{6.75}{\percent} إذا كان المصهور
مستخلصًا تجزيئيًا ومجمَّعًا).
\end{solution}
